\chapter{Signals, Modulation, and Receiver Architecture}
\label{ch:signals}
Radio carries information by changing a sinusoid, then recovers that information
from a selected portion of the spectrum.  This chapter starts with products of
sines, connects them to bandwidth, and follows the signal through a receiver.
It uses the mixer multiplication of \cref{sec:mixers}, the filter models of
\cref{ch:filters}, and the sampling view of \cref{ch:dsp}.

\section{Amplitude Modulation Makes Two Sidebands}
\label{sec:am}
Let a single-tone message be \(\cos\omega_m t\), a carrier be
\(\cos\omega_c t\), and \(m\) be the modulation index.  Conventional AM is
\[
s(t)=A_c[1+m\cos\omega_m t]\cos\omega_c t.
\]
Using \(\cos a\cos b=[\cos(a+b)+\cos(a-b)]/2\) gives
\[
\boxed{\;s(t)=A_c\cos\omega_ct+
\frac{mA_c}{2}\cos(\omega_c+\omega_m)t+
\frac{mA_c}{2}\cos(\omega_c-\omega_m)t.\;}
\]
One time-domain multiplication has become a carrier and equal upper and lower
sidebands.  A message whose highest audio frequency is \(f_m\) therefore
occupies \(2f_m\) of RF bandwidth.  This is Fourier analysis in a practical
form: multiplication shifts a spectrum, while a linear filter independently
scales each shifted component (\cref{sec:fourier}).

With a fixed load, power is proportional to squared amplitude.  The carrier
power is \(P_c\); each sideband has power \(m^2P_c/4\), so
\[
\boxed{\quad P_{\rm AM}=P_c\left(1+\frac{m^2}{2}\right).\quad}
\]
At \(m=1\), the sidebands together carry one third of the transmitted power.
Overmodulation, \(m>1\), makes the envelope cross zero and an envelope detector
can no longer recover the original message without distortion.

\section{SSB, Frequency Modulation, and Phase Modulation}
\label{sec:anglemod}
An SSB transmitter removes the carrier and one of AM's sidebands.  It needs only
the message bandwidth rather than twice the message bandwidth, and it does not
spend transmitter power on an uninformative carrier.  A receiver restores a
local carrier with a product detector; the mixer sum and difference terms of
\cref{sec:mixers} bring the selected sideband back to audio.

Frequency and phase modulation keep the carrier amplitude nearly constant and
put information into its angle.  A single-tone FM signal is
\[
s(t)=A_c\cos[\omega_ct+\beta\sin\omega_mt],\qquad
\beta=\frac{\Delta f}{f_m},
\]
where \(\Delta f\) is peak frequency deviation and \(\beta\) is the modulation
index.  Differentiating the phase shows the instantaneous frequency:
\[
f_i(t)=f_c+\Delta f\cos\omega_mt.
\]
FM has infinitely many mathematical sidebands, with amplitudes set by Bessel
functions; for engineering channel planning, Carson's rule keeps the important
ones:
\[
\boxed{\quad B\approx2(\Delta f+f_m).\quad}
\]
It is an approximation, not a hard spectral wall.  Phase modulation instead
makes phase directly proportional to the message.  Because frequency is the
time derivative of phase, an integrator before a phase modulator produces FM,
and a differentiator before an FM modulator produces PM.

\begin{exambox}
AM's bandwidth comes from two translated copies of the message.  SSB retains one.
FM's deviation is a frequency excursion, not its channel width: audio bandwidth
and deviation both enter Carson's rule.
\end{exambox}

\section{Why a Superheterodyne Uses an Intermediate Frequency}
\label{sec:receiverarchitecture}
A receiver wants to select one desired RF signal, amplify it without overload,
and recover its information.  In a superheterodyne the mixer translates the
desired RF to a fixed intermediate frequency (IF):
\[
\boxed{\quad f_{\rm IF}=|f_{\rm RF}-f_{\rm LO}|.\quad}
\]
Fixed-frequency filters can then supply repeatable selectivity and gain.  The
price is an image: a second RF frequency also differs from the local oscillator
by \(f_{\rm IF}\).  For high-side injection,
\[
f_{\rm LO}=f_{\rm RF}+f_{\rm IF},\qquad
\boxed{\quad f_{\rm image}=f_{\rm RF}+2f_{\rm IF}.\quad}
\]
An RF preselector must reject that image before the mixer, because after mixing
the wanted signal and image occupy the same IF and no later filter can tell them
apart.

The resulting signal path is a sequence of models already developed in this book:
\[
\text{antenna}\ \longrightarrow\ \text{preselector/LNA}\ \longrightarrow\
\text{mixer}\ \longrightarrow\ \text{IF filter and gain}\ \longrightarrow\
\text{detector or ADC}.
\]
The preselector trades bandwidth against image rejection and strong-signal
protection.  The first low-noise stage matters because of Friis's formula
(\cref{sec:friis}); the IF filter sets both adjacent-signal rejection and much of
the noise bandwidth (\cref{sec:noisebw}).  A direct-conversion or software-defined
receiver simply moves the translation and filtering into quadrature baseband and
digital processing; sampling still replicates spectra exactly as in \cref{ch:dsp}.

\begin{workedbox}[: one AM signal through a receiver]
An AM station at \SI{14.200}{\mega\hertz} carries a \SI{1}{\kilo\hertz} tone.
Its spectral lines are at \SI{14.199}{\mega\hertz}, \SI{14.200}{\mega\hertz},
and \SI{14.201}{\mega\hertz}; the occupied width is \SI{2}{\kilo\hertz}.
Choose \(f_{\rm IF}=\SI{9}{\mega\hertz}\) and high-side injection.  The local
oscillator is \SI{23.200}{\mega\hertz}, and a station at
\(14.200+2\times9=\SI{32.200}{\mega\hertz}\) produces the same IF.  A filter
after the mixer cannot remove it, so the RF preselector must do that work.
\end{workedbox}

\begin{mistakebox}
The IF is not an extra signal produced by the transmitter.  It is a frequency
chosen by the receiver designer so that its filters, gain, and detectors work at
one convenient frequency while the tuning occurs at RF.
\end{mistakebox}
